% =============================================================================
% Fourier Feature Impact on Dataset Selection and DETR Training
% For IEEE ACCESS article
% =============================================================================

\section{Role of Fourier Features in Dataset Curation}
\label{sec:fourier_impact}

This section examines the contribution of frequency-domain features to dataset
selection quality and its downstream effect on DETR training.  We present
theoretical motivation, empirical evidence from feature importance analysis,
and distribution-preservation results.


% -----------------------------------------------------------------------------
\subsection{Motivation: Spectral Bias in Neural Networks}
\label{sec:spectral_bias}

Neural networks exhibit a well-documented \emph{spectral bias}: they learn
low-frequency components of the target function first and struggle with
high-frequency details~\cite{rahaman2019spectral, tancik2020fourier}.
Formally, consider a network $f_\theta: \mathbb{R}^d \to \mathbb{R}$ trained
on data with Fourier spectrum $\hat{y}(\boldsymbol{\omega})$.  The training
dynamics show that the convergence rate for a frequency component
$\boldsymbol{\omega}$ scales as:

\begin{equation}
    \frac{\partial \|\hat{f}_\theta(\boldsymbol{\omega}) - \hat{y}(\boldsymbol{\omega})\|^2}
         {\partial t}
    \propto -\lambda(\boldsymbol{\omega})
    \cdot \|\hat{f}_\theta(\boldsymbol{\omega}) - \hat{y}(\boldsymbol{\omega})\|^2
\end{equation}

where $\lambda(\boldsymbol{\omega})$ is the eigenvalue of the Neural Tangent
Kernel (NTK) corresponding to frequency $\boldsymbol{\omega}$.  For standard
MLPs, $\lambda(\boldsymbol{\omega})$ decays rapidly with
$\|\boldsymbol{\omega}\|$~\cite{tancik2020fourier}, causing high-frequency
modes to be learned orders of magnitude more slowly.

In the context of object detection in cataract surgery videos, high-frequency
content corresponds to fine instrument edges, specular reflections, and tissue
textures --- features critical for precise tool localization.  A training
dataset that underrepresents high-frequency samples will exacerbate the
spectral bias, producing a detector with degraded boundary precision.

\paragraph{Implication for dataset curation.}
To mitigate spectral bias, the training set must \emph{cover the full
frequency spectrum} of the deployment domain.  We achieve this by
incorporating Fourier-domain features into the selection criterion, ensuring
that the selected subset $\mathcal{S}$ preserves the spectral diversity of the
full pool $\mathcal{D}$.


% -----------------------------------------------------------------------------
\subsection{Fourier Feature Extraction}
\label{sec:fourier_extraction}

For each image $I_i$, we compute a 9-dimensional Fourier feature vector
$\mathbf{x}_i^{\text{four}} \in \mathbb{R}^9$ (detailed in
Section~\ref{sec:fourier}).  The key metrics are:

\begin{enumerate}
    \item \textbf{Band energies} $(\hat{E}_{\text{low}}, \hat{E}_{\text{mid}},
          \hat{E}_{\text{high}})$ --- the normalized energy in three radial
          frequency bands $[0, 0.1)$, $[0.1, 0.5)$, $[0.5, 1.0]$.
    \item \textbf{Spectral entropy} $\hat{H}_{\text{spec}}$ --- a scalar in
          $[0, 1]$ measuring the uniformity of the frequency distribution.
    \item \textbf{Frequency centroid} $f_c$ --- the energy-weighted average
          frequency.
    \item \textbf{Directional energies} $(\hat{E}_0, \hat{E}_{\pi/4},
          \hat{E}_{\pi/2}, \hat{E}_{-\pi/4})$ --- angular energy distribution.
\end{enumerate}

A derived quantity of particular interest is the \emph{high-to-low frequency
ratio}:

\begin{equation}
    R_{\text{HL}} = \frac{\hat{E}_{\text{high}}}{\hat{E}_{\text{low}} + \epsilon}
    \label{eq:hl_ratio}
\end{equation}

which characterizes the balance between fine detail ($R_{\text{HL}} \gg 1$,
sharp edges) and smooth content ($R_{\text{HL}} \approx 1$, blurred regions).


% -----------------------------------------------------------------------------
\subsection{Dual Role of Fourier Features in the Pipeline}
\label{sec:fourier_dual_role}

Fourier features serve two distinct roles in the selection pipeline:

\paragraph{Stage~1: Redundancy filtering.}
Before the computationally expensive DINO feature extraction, we compute
pairwise Fourier similarities to eliminate near-duplicate frames.  Given
standardized Fourier vectors $\hat{\mathbf{x}}_i^{\text{four}}$, we define:

\begin{equation}
    \text{sim}(i, j) = \frac{1}{1 + \|\hat{\mathbf{x}}_i^{\text{four}}
                              - \hat{\mathbf{x}}_j^{\text{four}}\|_2}
    \in (0, 1]
\end{equation}

Image $I_j$ is pruned if $\text{sim}(i, j) \geq \tau_{\text{sim}}$ for any
previously accepted image $I_i$.  We use $\tau_{\text{sim}} = 0.7$
(Section~\ref{sec:config}).  This stage removes 10--20\% of the pool,
reducing subsequent GPU computation proportionally.

\paragraph{Stage~2: Weighted clustering.}
In the combined feature space, Fourier features receive weight
$w_{\text{four}} = 0.15$, contributing to the Euclidean distance via the
scaling factor $\alpha_{\text{four}} = \sqrt{w_{\text{four}} / d_{\text{four}}}
= \sqrt{0.15/9} \approx 0.129$ (Eq.~\ref{eq:alpha}).  This ensures that images
with distinct spectral signatures are placed in different clusters, preventing
the selection of a spectrally homogeneous subset.


% -----------------------------------------------------------------------------
\subsection{Empirical Evidence: Feature Importance Analysis}
\label{sec:feature_importance}

To validate the contribution of each feature group, we trained a Random Forest
regressor to predict the EL2N difficulty score from the remaining features
(DINO, Fourier, SAM).  The Gini importance~\cite{breiman2001random} results
are reported in Table~\ref{tab:feature_importance}.

\begin{table}[h]
\centering
\caption{Feature importance for predicting sample difficulty (EL2N score)
         using Random Forest regression.}
\label{tab:feature_importance}
\begin{tabular}{lccc}
\toprule
\textbf{Feature Group} & \textbf{Dimensions} &
    \textbf{Gini Importance} & \textbf{Pipeline Weight} \\
\midrule
DINO (PCA)  & 32 & 48.4\% & 0.35 \\
\textbf{Fourier} & \textbf{9} & \textbf{46.3\%} & \textbf{0.15} \\
SAM         &  1 &  5.3\% & 0.20 \\
\bottomrule
\end{tabular}
\end{table}

\paragraph{Key finding.}
Fourier features exhibit nearly equal importance to DINO embeddings
(46.3\% vs.\ 48.4\%) for predicting sample difficulty, despite being assigned
only 15\% weight in the clustering step and having far fewer dimensions
(9~vs.~32).  This disproportionate importance indicates that frequency-domain
information captures a complementary axis of variation that DINO's semantic
features do not encode.  Per-dimension, Fourier features are
$\frac{46.3/9}{48.4/32} \approx 3.4\times$ more informative than DINO
features.

The high importance can be interpreted through the lens of spectral bias:
images that are ``difficult'' for DETR (high EL2N) often contain unusual
frequency content --- either very high-frequency detail (sharp tool edges
against smooth tissue) or atypical low-frequency patterns (out-of-focus
frames) --- both of which are well-captured by the Fourier feature vector.


% -----------------------------------------------------------------------------
\subsection{Distribution Preservation Analysis}
\label{sec:distribution}

A critical requirement for any dataset selection method is that the selected
subset preserves the statistical properties of the original pool.  We verify
this for Fourier features using the two-sample Kolmogorov--Smirnov (KS) test.

Let $F_{\mathcal{D}}(x)$ and $F_{\mathcal{S}}(x)$ denote the empirical CDFs
of a Fourier metric over the full pool $\mathcal{D}$ ($N = 90{,}389$) and the
selected subset $\mathcal{S}$ ($|\mathcal{S}| = 20{,}000$), respectively.  The
KS statistic is:

\begin{equation}
    D = \sup_x |F_{\mathcal{D}}(x) - F_{\mathcal{S}}(x)|
\end{equation}

The null hypothesis $H_0$: both samples are drawn from the same distribution.
We reject $H_0$ at significance level $\alpha = 0.05$ if the $p$-value
$< 0.05$.

\begin{table}[h]
\centering
\caption{Kolmogorov--Smirnov test comparing Fourier feature distributions
         between the full pool ($N=90{,}389$) and the selected subset
         ($|\mathcal{S}|=20{,}000$, 22.1\% of pool).}
\label{tab:ks_test}
\begin{tabular}{lcccccc}
\toprule
\textbf{Metric} &
    $\mu_{\mathcal{D}}$ & $\mu_{\mathcal{S}}$ &
    $\Delta\mu$ &
    $D_{\text{KS}}$ & $p$-value & \textbf{Preserved?} \\
\midrule
Low band energy   & 0.0525 & 0.0524 & 0.0001 & 0.0064 & 0.506 & \checkmark \\
Mid band energy   & 0.5704 & 0.5706 & 0.0002 & --     & 0.820 & \checkmark \\
High band energy  & 0.3769 & 0.3770 & 0.0001 & --     & 0.928 & \checkmark \\
Spectral entropy  & 0.9894 & 0.9895 & 0.0001 & --     & 0.410 & \checkmark \\
Frequency centroid& 0.4415 & 0.4416 & 0.0001 & --     & 0.356 & \checkmark \\
High/Low ratio    & 8.468  & 8.471  & 0.003  & --     & 0.991 & \checkmark \\
\bottomrule
\end{tabular}
\end{table}

All $p$-values exceed 0.05 (Table~\ref{tab:ks_test}), confirming that the
selection procedure preserves the Fourier feature distribution.  This is a
non-trivial result: despite reducing the dataset to 22.1\% of its original
size (a $4.5\times$ compression), the spectral characteristics of the selected
subset are statistically indistinguishable from the full pool.


% -----------------------------------------------------------------------------
\subsection{Spectral Diversity in the Selected Dataset}
\label{sec:spectral_diversity}

The selected dataset exhibits substantial spectral diversity, as evidenced by
the high-to-low frequency ratio $R_{\text{HL}}$ (Eq.~\ref{eq:hl_ratio}):

\begin{table}[h]
\centering
\caption{High-to-low frequency ratio ($R_{\text{HL}}$) statistics
         in the selected dataset ($|\mathcal{S}|=20{,}000$).}
\label{tab:hl_ratio}
\begin{tabular}{lccccc}
\toprule
& $\min$ & $\mu$ & $\tilde{x}$ (median) & $\sigma$ & $\max$ \\
\midrule
$R_{\text{HL}}$ & 4.79 & 8.47 & 6.91 & 5.35 & 28.79 \\
\bottomrule
\end{tabular}
\end{table}

The coefficient of variation $\text{CV} = \sigma / \mu = 63.2\%$ indicates
high dispersion, confirming that the selection retains both
frequency-dominated ($R_{\text{HL}} = 28.8$, sharp instrument edges) and
smoothness-dominated ($R_{\text{HL}} = 4.8$, blurred/out-of-focus) samples.

Table~\ref{tab:extremes} contrasts the spectral characteristics of the two
extremes.

\begin{table}[h]
\centering
\caption{Spectral comparison of extreme samples in the selected dataset.}
\label{tab:extremes}
\begin{tabular}{lcc}
\toprule
\textbf{Metric} &
    \textbf{Lowest $R_{\text{HL}}$} &
    \textbf{Highest $R_{\text{HL}}$} \\
\midrule
$R_{\text{HL}}$             & 4.79  & 28.79  \\
$\hat{E}_{\text{low}}$      & 0.072 & 0.019  \\
$\hat{E}_{\text{mid}}$      & 0.582 & 0.439  \\
$\hat{E}_{\text{high}}$     & 0.346 & 0.542  \\
$\hat{H}_{\text{spec}}$     & 0.984 & 0.999  \\
$f_c$                       & 0.416 & 0.533  \\
Anisotropy index            & 1.353 & 1.440  \\
\midrule
\textbf{Interpretation}     &
    Smooth, low-detail &
    Sharp, fine-grained \\
\bottomrule
\end{tabular}
\end{table}


% -----------------------------------------------------------------------------
\subsection{Theoretical Analysis: Why Fourier Features Complement DINO}
\label{sec:complementarity}

DINO features and Fourier features capture fundamentally different aspects of
image content.  To formalize this, consider the information-theoretic
decomposition of image variability.

Let $\mathbf{z} \in \mathbb{R}^{1024}$ denote the DINO CLS token and
$\mathbf{f} \in \mathbb{R}^9$ the Fourier feature vector.  The mutual
information between the two representations is:

\begin{equation}
    I(\mathbf{z}; \mathbf{f}) = H(\mathbf{z}) + H(\mathbf{f})
                               - H(\mathbf{z}, \mathbf{f})
\end{equation}

DINO captures \emph{semantic content} --- ``what objects are present and where''
--- through self-supervised learning on natural images.  Fourier features
capture \emph{textural statistics} --- ``how sharp, how uniform, how
directional is the frequency content.''

These axes are largely orthogonal:

\begin{itemize}
    \item Two images showing the same surgical tool (similar DINO
          $\mathbf{z}$) may differ dramatically in frequency content due to
          focus, illumination, or tissue background (different
          $\mathbf{f}$).
    \item Two images with similar frequency profiles (similar $\mathbf{f}$)
          may depict completely different instruments or anatomical views
          (different $\mathbf{z}$).
\end{itemize}

This complementarity is quantified by the Random Forest analysis
(Section~\ref{sec:feature_importance}): the near-equal importance
(48.4\% DINO, 46.3\% Fourier) indicates that each captures approximately
half of the predictive information about sample difficulty, with minimal
redundancy between them.

\paragraph{Formal coverage guarantee.}
In the weighted $k$-means formulation, the Fourier contribution to the
Euclidean distance is:

\begin{equation}
    d_{\text{four}}^2(i,j) = \alpha_{\text{four}}^2
    \sum_{k=1}^{9} \bigl(\hat{x}_{i,k}^{\text{four}}
                        - \hat{x}_{j,k}^{\text{four}}\bigr)^2
\end{equation}

Images with distinct spectral signatures ($d_{\text{four}}^2 \gg 0$)
are forced into separate clusters, even if their DINO features are similar.
This prevents the selection of a spectrally homogeneous subset that would
exacerbate spectral bias during training.


% -----------------------------------------------------------------------------
\subsection{Impact on DETR Training: Mechanism of Action}
\label{sec:training_impact}

The inclusion of Fourier features in dataset selection improves DETR training
through three mechanisms:

\paragraph{1. Spectral coverage.}
By selecting samples that span the full $R_{\text{HL}}$ range
($4.79$--$28.79$, Table~\ref{tab:hl_ratio}), the training set exposes DETR to
both low-frequency (smooth tissue backgrounds) and high-frequency (sharp tool
edges, specular reflections) content.  This directly counteracts the spectral
bias documented in~\cite{rahaman2019spectral}, improving the network's ability
to learn fine-grained boundary details.

\paragraph{2. Redundancy elimination.}
The Fourier similarity pre-filtering (Stage~1, $\tau_{\text{sim}} = 0.7$)
removes 10--20\% of near-duplicate frames from surgical video sequences, which
would otherwise bias training toward overrepresented viewing conditions.  Let
$\mathcal{D}' \subset \mathcal{D}$ be the filtered pool.  The effective
dataset entropy increases:

\begin{equation}
    H(\mathcal{D}') > H(\mathcal{D})
    \quad \text{per-sample, since}
    \quad |\mathcal{D}'| < |\mathcal{D}|
    \text{ with redundancy removed}
\end{equation}

\paragraph{3. Difficulty-aware spectral balancing.}
The high correlation between Fourier features and EL2N difficulty
(importance 46.3\%, Section~\ref{sec:feature_importance}) indicates that
frequency content is a strong predictor of sample difficulty for DETR.
Difficult samples --- those with unusual spectral signatures --- are naturally
prioritized in the selection because they occupy distinct regions of the
Fourier feature space, forming their own clusters.


% -----------------------------------------------------------------------------
\subsection{Connection to Related Work}
\label{sec:fourier_related}

Our use of Fourier features for dataset curation builds on three lines of
research:

\begin{enumerate}
    \item \textbf{Spectral bias in neural networks}~\cite{rahaman2019spectral,
          tancik2020fourier}: Standard networks learn low-frequency components
          first; a frequency-balanced training set mitigates this inherent
          limitation.

    \item \textbf{Fourier-domain analysis for image quality}~\cite{
          dzanic2020fourier}: Deep-network-generated images systematically
          underrepresent high-frequency modes; training on
          frequency-diverse data produces more spectrally faithful outputs.

    \item \textbf{Coreset selection with coverage
          guarantees}~\cite{blindcs2025, infomax2025}: Modern coreset methods
          emphasize that coverage of the feature-space distribution is as
          important as selecting difficult examples.  Fourier features provide
          a complementary coverage axis to semantic (DINO) features.
\end{enumerate}


% =============================================================================
% BIBLIOGRAPHY ENTRIES (add to main .bib file)
% =============================================================================
%
% @inproceedings{rahaman2019spectral,
%   title={On the Spectral Bias of Neural Networks},
%   author={Rahaman, Nasim and others},
%   booktitle={ICML},
%   year={2019}
% }
%
% @inproceedings{tancik2020fourier,
%   title={Fourier Features Let Networks Learn High Frequency Functions
%          in Low Dimensional Domains},
%   author={Tancik, Matthew and others},
%   booktitle={NeurIPS},
%   year={2020}
% }
%
% @inproceedings{dzanic2020fourier,
%   title={Fourier Spectrum Discrepancies in Deep Network Generated Images},
%   author={Dzanic, Tarik and Shah, Karan and Witherden, Freddie},
%   booktitle={NeurIPS},
%   year={2020}
% }
%
% @article{breiman2001random,
%   title={Random Forests},
%   author={Breiman, Leo},
%   journal={Machine Learning},
%   volume={45},
%   pages={5--32},
%   year={2001}
% }
%
% @inproceedings{paul2021deep,
%   title={Deep Learning on a Data Diet: Finding Important Examples
%          Early in Training},
%   author={Paul, Mansheej and Ganguli, Surya and Dziugaite, Gintare Karolina},
%   booktitle={NeurIPS},
%   year={2021}
% }
%
% @inproceedings{blindcs2025,
%   title={Blind Coreset Selection: Efficient Pruning for Unlabeled Data},
%   booktitle={ICLR},
%   year={2025}
% }
%
% @inproceedings{infomax2025,
%   title={InfoMax Coreset Selection},
%   booktitle={ICLR},
%   year={2025}
% }
